\pdfoutput=1
\documentclass[11pt]{article}
\usepackage[review]{acl2023}
\usepackage{times}
\usepackage{latexsym}
\usepackage[T1]{fontenc}
\usepackage[utf8]{inputenc}
\usepackage{microtype}
\usepackage{inconsolata}
\usepackage{booktabs}
\usepackage{xcolor}

\newcommand{\pending}[1]{\textcolor{red}{\textbf{[PENDING: #1]}}}

\title{LLM-Primary Gold Annotation for Cyber Threat Intelligence Named Entity Recognition: Honest Evaluation of Token-Classification Architectures}

\author{Daniel Dye \\
  NLP Capstone, MS in Artificial Intelligence \\
  University of South Florida \\
  \texttt{[email]} }

\begin{document}
\maketitle

\begin{abstract}
Cyber threat intelligence (CTI) reports contain entity types with no
general-domain analogue --- ATT\&CK technique identifiers, file hashes, TLP
markings, threat actor aliases --- that off-the-shelf NER models handle
poorly. Prior work in this space, including an earlier iteration of this
project, evaluates against \emph{silver} labels produced by the same class
of dictionary/rule annotator being improved upon, which we show is
circular: a checkpoint that scores 0.74 F1 against dictionary-derived
silver gold collapses to 0.35 against independently-produced gold, because
the silver labels share the model's blind spots rather than measuring
against them. We instead build gold directly with a frontier LLM (Claude
Sonnet), which reads full advisory text and proposes exhaustive entity
annotations under per-label span validators, followed by human curation on
the evaluation split. We fine-tune three token-classification encoders ---
RoBERTa-base, RoBERTa-large \citep{liu2019roberta}, and SecureBERT
\citep{aghaei2022securebert}, a RoBERTa-architecture model domain-pretrained
on security text --- on this gold and evaluate all three with an identical
document-level harness. SecureBERT is the best performer (0.71 strict
micro-F1), ahead of RoBERTa-base (0.66) and RoBERTa-large (0.65), showing
domain pretraining matters more than parameter count on this task. All
three models are near-perfect on dictionary-shaped structured types
(attack technique id, CVE identifier, TLP classification, email address)
and substantially weaker on semantic types with fewer, more varied training
examples (threat actor group, hacking tool). \pending{a GLiNER checkpoint
retrained on the identical LLM gold, and the corresponding architecture
comparison, is in progress and will be added once training completes}.
\end{abstract}

\section{Introduction}

Cyber threat intelligence reports --- government advisories, malware
analyses, incident write-ups --- are the primary vehicle by which defenders
share knowledge, yet extracting structured information from them remains
largely manual. Named entity recognition over CTI text must handle entity
types with no general-domain analogue: ATT\&CK technique identifiers, file
hashes, TLP markings, YARA rule names, threat actor aliases. Human
annotation at scale is expensive in this domain because it requires
security expertise, which motivated an earlier iteration of this project to
avoid it entirely: build a \emph{silver}-labeled corpus from agreement
between a dictionary/regex annotator (STIX indicators, MITRE ATT\&CK
lookups) and zero-shot GLiNER \citep{zaratiana2024gliner} predictions, and
evaluate against held-out silver labels produced by the same dictionary
annotator.

That evaluation is circular, and the gap it hides is large. A blind human
audit of predicted spans found the dictionary annotator's own gold-recall
was materially incomplete --- it misses real entities the dictionary simply
never encoded --- so a model that learns to reproduce the dictionary's
blind spots is rewarded, not penalized, by silver-label scoring. We confirm
this quantitatively: rescoring an existing dictionary-silver-trained
checkpoint against independently-produced gold drops its score from 0.74 to
0.35 F1 (\S\ref{sec:results}). The model did not get worse; the measurement
got honest.

This motivates the pipeline described here. Rather than avoiding human
involvement entirely, we move it to where it is cheapest to apply
correctly: a frontier LLM (Claude Sonnet) reads each advisory's full text
and proposes exhaustive gold-standard entity annotations, subject to
per-label span-format validators at generation time; a human then curates
the evaluation split (accepting, rejecting, or relabeling proposals), while
the larger training split remains LLM-only. This is still not free of
human effort, but it replaces exhaustive first-pass annotation with
targeted review --- and, critically, it produces evaluation gold that is
\emph{independent} of any dictionary or rule system a model might learn to
imitate.

We ask four research questions:

\begin{description}
  \item[RQ1] How large is the gap between circular (dictionary-silver)
    evaluation and honest (LLM-primary, human-curated) evaluation for a
    model trained on the old weak-supervision pipeline?
  \item[RQ2] Once both training and evaluation labels come from the honest
    LLM-gold pipeline, does model scale (RoBERTa-base vs.\ RoBERTa-large)
    or domain-adaptive pretraining (SecureBERT) matter more for
    token-classification CTI NER?
  \item[RQ3] Which entity types remain hard even after retraining on
    honest gold, and does the pattern track training-example frequency,
    structured-vs-semantic label type, or both?
  \item[RQ4] How does this fine-tuned token-classification fleet compare
    to a zero-shot span-matching architecture (GLiNER) retrained on the
    identical honest gold, and to a frontier LLM run zero-shot?
\end{description}

RQ1--RQ3 are answered here with data trained and scored on the current
LLM-gold pipeline. RQ4 requires a GLiNER checkpoint retrained on the
identical gold, which is in progress and reported as \pending{} throughout;
the frontier-LLM comparison is deferred to future work for the same
reason --- the Gemini numbers from the prior weak-supervision iteration of
this project are not comparable to the current gold and are not reported
here.

\section{Related Work}

GLiNER \citep{zaratiana2024gliner} treats NER as matching natural-language
entity-type descriptions against learned span embeddings, enabling
zero-shot extraction of arbitrary types; its schema-driven successor
GLiNER2 \citep{zaratiana2025gliner2} unifies NER with classification and
structured extraction in one model. Distant supervision for NER
\citep{devlin2019bert} was formalized for open-domain text by BOND
\citep{liang2020bond}, which fine-tunes a token classifier on
knowledge-base-derived labels then self-trains through a teacher-student
loop to denoise them --- the same failure modes we target with ensemble
agreement (incomplete and noisy weak labels) motivate BOND's two-stage
design.

Cybersecurity-specific NER corpora suffer acutely from schema fragmentation:
\citet{schmidt2025labelunification} show that naively unifying four
CTI datasets under a coarse shared label set \emph{hurts} performance
relative to single-dataset training, while \citet{echchammakhy2025cyberner}
demonstrate that principled, STIX-anchored harmonization recovers a 30\%
relative F1 gain over the same naive union --- evidence that \emph{how}
labels are unified matters as much as how much data is unified; our
LLM-primary annotation pipeline sidesteps unification entirely by defining
the label schema once and generating gold directly for it, rather than
merging pre-existing dictionary schemas. \citet{deka2024attacker} extend
the entity-type inventory to attribution-specific concepts (impact,
course-of-action) with deliberately long, sentence-level spans, and find
instruction-fine-tuned LLMs outperform encoder taggers on their
benchmark, motivating our use of an LLM as the \emph{annotator} rather
than the deployed extractor. \citet{alam2024ctibench} benchmark LLMs
directly on CTI tasks including entity extraction, providing the
frontier-model baseline methodology RQ4's future-work extension will
adapt.

LLM-as-annotator is itself an active line of work distinct from
LLM-as-extractor: rather than running a frontier model zero-shot at
inference time, we use it once, offline, to produce training and
evaluation labels for smaller, cheaper, deployable encoders --- a
distillation-flavored setup, though unlike classic distillation the
teacher's output here is explicitly checked (span-format validators,
human curation) rather than trusted directly, because the failure mode we
are guarding against is exactly the one BOND and the dictionary-silver
approach share: an annotator's systematic blind spots being learned
uncritically by the student.

\section{Data and Method}
\label{sec:data-method}

\subsection{Pipeline}

We scrape approximately 5{,}000 CISA cybersecurity advisories; each is
converted from PDF to markdown via docling, then stripped of layout
artifacts and normalized (page-furniture dedup, pagination removal,
split-hash rejoining) to a clean, model-facing text version. For 82 of
these advisories we generate gold entity annotations across 16 CTI-specific
labels\footnote{affected software product, attack technique id, attack
technique name, cve identifier, cwe identifier, email address, file hash,
file name, file path, hacking tool, ip address, malware family, tactic,
threat actor group, tlp classification, yara rule name.} directly with a
frontier LLM (Claude Sonnet): the model reads each advisory's full clean
text and proposes exhaustive entity spans, verbatim-matched back to the
source text and rejected if they do not match exactly, and further gated by
per-label span-format validators (hash length and charset, CVE/CWE/ATT\&CK
technique-ID grammar, IP/email format including defanged forms, TLP
enumeration) so structurally malformed proposals never enter the gold set.
This is an annotation-time generation gate, not a post-hoc filter.

The 66 training-split advisories receive LLM annotation only. The 16
validation/test advisories additionally go through human curation via a
purpose-built review tool: a human accepts, rejects, or relabels each LLM
proposal and may add entities the LLM missed, including an \texttt{UNK}
category (not used for training or evaluation directly) for real entities
that fit no current label, which we use to look for missing label types.
Evaluation gold is therefore human-verified; training gold is not, which we
report as a limitation (\S\ref{sec:limitations}) but consider acceptable
because the LLM-primary approach's failure mode --- proposing an entity a
human would reject --- is structurally different from the dictionary
approach's failure mode of \emph{never proposing} entire classes of entity
it has no rule for, which is what a human-verified evaluation set is
positioned to catch either way.

The training split contains 6{,}701 gold entities across 66 advisories;
the evaluation split (8 validation, 8 test advisories) contains 1{,}995
gold spans in the test set alone once chunked for scoring
(\S\ref{sec:evaluation}). This is roughly 3.5$\times$ the raw entity volume
of the prior dictionary-silver pipeline's test set on the same page count,
because dictionary/regex annotation structurally omits several semantic
label types (e.g.\ \emph{affected software product}, \emph{hacking tool})
that an LLM reader identifies from context.

\subsection{Split}

We split by advisory \emph{document}, not by chunk, before chunking for
training --- 66 advisories train / 8 validation / 8 test --- so that
overlapping text windows from the same document cannot appear on both
sides of the split. This was a deliberate fix carried over from the prior
pipeline iteration: chunk-level splitting leaks document context (shared
boilerplate, repeated entity mentions) across the train/test boundary and
inflates apparent performance.

\subsection{Models}

\paragraph{RoBERTa-base and RoBERTa-large \citep{liu2019roberta}.} Standard
BIO token classifiers fine-tuned on the LLM-gold training split. Because
gold spans are word-level and RoBERTa's tokenizer is subword-based, we tag
the first wordpiece of each word and mask continuation wordpieces from the
loss. Corpus chunks vary widely in length, and long or hash-heavy chunks
fragment into more wordpieces than a single window admits under naive
truncation; we instead tokenize with overlapping sliding windows (a
510-wordpiece budget per window, 50-word overlap between windows) so every
word is covered by at least one window, and take each word's predicted tag
from the window in which it sits most centrally at inference time.
RoBERTa-base trains for 8 epochs at LR 2e-5, effective batch size 16;
RoBERTa-large under the identical recipe with effective batch size 32
(gradient accumulation, to fit the larger model). Length-grouped batching
reduces padding waste across the wide chunk-length range.

\paragraph{SecureBERT \citep{aghaei2022securebert}.} The same
RoBERTa architecture and fine-tuning recipe as RoBERTa-base (8 epochs, LR
2e-5, effective batch 16), but initialized from a checkpoint further
pretrained on a security-text corpus rather than general-domain text. Since
architecture and fine-tuning recipe are held identical to RoBERTa-base, any
difference between the two isolates the effect of domain-adaptive
pretraining from the effect of the fine-tuning data or procedure.

\paragraph{GLiNER (pending).} \texttt{urchade/gliner\_medium-v2.1}
\citep{zaratiana2024gliner}, zero-shot span-matching, is being fine-tuned
on the identical LLM-gold training split at the time of writing. We report
comparisons to it once training completes (\S\ref{sec:results},
\S\ref{sec:limitations}).

\paragraph{Frontier LLM (future work).} A zero-shot frontier-LLM baseline
on the current gold, in the style of \citet{alam2024ctibench}, is deferred
to future work; the Gemini numbers from the prior pipeline iteration are
scored against dictionary-silver gold and are not comparable here.

\subsection{Evaluation}
\label{sec:evaluation}

All models are scored with the same document-level harness
(\texttt{ner.evaluate}): predictions from overlapping chunks are merged
back to document-global token indices before matching, so entities in
chunk-overlap regions are not double-counted. We report strict (exact
span-and-label match) and relaxed (label match plus span overlap)
micro/macro precision, recall, and F1, on the 8-document test split, using
the identical scoring implementation across models so the comparison in
Section~\ref{sec:results} is not an artifact of different matching rules.

\section{Results}
\label{sec:results}

\subsection{RQ1: Circular vs.\ Honest Evaluation}

Table~\ref{tab:circular-vs-honest} rescores the prior pipeline's checkpoints
--- trained on dictionary-silver labels, previously reported at 0.74 and
0.76 F1 --- against the current, independently-produced LLM+human gold, on
the same 8 test advisories (1{,}995 gold spans) used throughout this paper.
Both token-classification checkpoints collapse: RoBERTa-base from 0.740 to
0.351, SecureBERT from 0.755 to 0.393, roughly a 47--48\% relative drop.
The model did not get worse; the measurement got honest --- these
checkpoints were never evaluated against labels their own dictionary
annotator didn't also produce, so silver scoring rewarded them for learning
the dictionary's blind spots along with its correct patterns. For
reference, the old GLiNER checkpoint (fine-tuned on the same dictionary
silver labels, not yet retrained on the current gold) drops by a smaller
margin, 0.666 to 0.540 (19\% relative). Whether that reflects a genuine
architectural advantage in span-matching generalization or is an artifact
of GLiNER's own training data being marginally less circular is exactly
the question RQ4 needs a GLiNER checkpoint retrained on \emph{this} gold to
answer, and is left pending.

\begin{table}[t]
\centering
\small
\begin{tabular}{lrrrr}
\toprule
Model & Circular F1 & Honest P & Honest R & Honest F1 \\
\midrule
RoBERTa-base & 0.740 & 0.86 & 0.22 & 0.351 \\
SecureBERT & 0.755 & 0.87 & 0.25 & 0.393 \\
GLiNER (v1, ref.) & 0.666 & 0.65 & 0.46 & 0.540 \\
\bottomrule
\end{tabular}
\caption{Checkpoints trained on the prior dictionary-silver pipeline, scored against the silver labels they were trained toward (``Circular'') vs.\ against independent LLM+human gold (``Honest''), strict span+label match, 8-document test split.}
\label{tab:circular-vs-honest}
\end{table}

\subsection{RQ2: Scale vs.\ Domain Pretraining Under Honest Gold}

Table~\ref{tab:architecture} reports the three encoders after retraining on
the LLM-gold training split and scoring against the same honest test gold.
SecureBERT is the best performer (0.712 strict / 0.758 relaxed micro-F1),
ahead of RoBERTa-base (0.659 / 0.683) and, notably, ahead of RoBERTa-large
(0.648 / 0.697) despite RoBERTa-large having roughly 3$\times$ the
parameters. Domain-adaptive pretraining recovers more honest-gold quality
here than raw scale does. RoBERTa-large's higher relaxed-F1 and macro-F1
alongside a \emph{lower} strict micro-F1 than RoBERTa-base is explained by
the per-label pattern in \S\ref{sec:rq3}: RoBERTa-large predicts the right
general region and type more consistently across labels (helping macro and
relaxed matching) but regresses sharply on the single highest-volume label,
\emph{malware family}, which drags support-weighted strict micro-F1 down.

Comparing back to Table~\ref{tab:circular-vs-honest}: retraining on honest
gold closes most of the gap the circular evaluation had hidden.
RoBERTa-base's honest score rises from 0.351 (old checkpoint) to 0.659
(retrained) --- most of the way back to its old, circular 0.740 --- and
SecureBERT's honest score (0.712) very nearly matches its old circular
number (0.755). Fixing \emph{both} the training labels and the evaluation
labels recovers almost all of the quality the circular number implied was
already there; fixing only the evaluation labels (Table~\ref{tab:circular-vs-honest})
is what made the models look like they had regressed.

\begin{table}[t]
\centering
\small
\begin{tabular}{lrrrr}
\toprule
Model & P & R & F1 & F1 (relaxed) \\
\midrule
RoBERTa-base & 0.71 & 0.61 & 0.659 & 0.683 \\
RoBERTa-large & 0.71 & 0.60 & 0.648 & 0.697 \\
SecureBERT & \textbf{0.73} & \textbf{0.70} & \textbf{0.712} & \textbf{0.758} \\
GLiNER & \multicolumn{4}{c}{\pending{retraining on identical gold in progress}} \\
\bottomrule
\end{tabular}
\caption{Token-classification architecture comparison, identical LLM-gold training and evaluation labels, identical evaluation harness. Micro P/R/F1, 8-document test split, 1{,}995 gold spans.}
\label{tab:architecture}
\end{table}

\subsection{RQ3: Which Entity Types Remain Hard}
\label{sec:rq3}

Table~\ref{tab:per-label} breaks the three retrained models down by label.
Structured, pattern-shaped labels are near-solved regardless of frequency:
\emph{attack technique id}, \emph{cve identifier}, \emph{tlp
classification}, and \emph{email address} all score 0.92--1.00 F1 across
all three models, and even the lowest-support label in the test set
(\emph{cwe identifier}, $n=4$) hits 1.00 on every model --- though at that
support level a single miss would swing the number by 0.25, so this
particular result should be read as indicative rather than as evidence
CWE identification is solved.

The weakest labels are not simply the rarest ones. \emph{malware family}
has the largest test-set support of any label ($n=438$) yet tops out at
0.65 F1 (SecureBERT); \emph{threat actor group} ($n=157$) and
\emph{hacking tool} ($n=139$) similarly lag well behind the structured
labels despite moderate-to-high support. This suggests frequency alone
does not explain the gap: these are open-vocabulary, context-dependent
labels (whether a string names a piece of malware, a named intrusion set,
or a legitimate admin tool like PowerShell often depends on surrounding
sentence context rather than surface form), and that semantic ambiguity
appears to be a harder problem for a token classifier than sparse
supervision is. RoBERTa-large's regression on \emph{malware family}
specifically (0.35 vs.\ 0.59 for RoBERTa-base and 0.65 for SecureBERT, on
identical training data) is the single largest per-label swing in the
table and the main driver of RoBERTa-large's lower strict micro-F1 in
Table~\ref{tab:architecture}.

\begin{table}[t]
\centering
\small
\begin{tabular}{lrrrr}
\toprule
Label & $n$ & Base & Large & SecBERT \\
\midrule
malware family             & 438 & 0.59 & 0.35 & 0.65 \\
attack technique id        & 297 & 0.86 & 0.91 & \textbf{0.99} \\
affected software product  & 159 & 0.48 & 0.57 & 0.61 \\
threat actor group         & 157 & 0.29 & 0.42 & 0.34 \\
attack technique name      & 144 & 0.57 & 0.53 & 0.57 \\
hacking tool                & 139 & 0.37 & 0.44 & 0.51 \\
ip address                 & 130 & 0.85 & 0.85 & 0.92 \\
tactic                     & 111 & 0.59 & 0.63 & 0.62 \\
file path                  & 108 & 0.84 & 0.83 & 0.84 \\
file hash                  &  93 & 0.70 & 0.65 & 0.64 \\
file name                  &  89 & 0.76 & 0.81 & 0.74 \\
cve identifier              &  55 & 0.92 & 0.96 & 0.94 \\
tlp classification         &  36 & 0.97 & 0.97 & 0.97 \\
email address              &  26 & 0.98 & 0.98 & 0.98 \\
yara rule name              &   9 & 0.62 & 0.71 & 0.50 \\
cwe identifier               &   4 & 1.00 & 1.00 & 1.00 \\
\midrule
\textbf{Micro} & 1995 & 0.659 & 0.648 & \textbf{0.712} \\
\bottomrule
\end{tabular}
\caption{Per-label strict F1, sorted by test-set support. Base = RoBERTa-base, Large = RoBERTa-large, SecBERT = SecureBERT.}
\label{tab:per-label}
\end{table}

\subsection{RQ4: Architecture and Frontier Comparisons (Pending)}
\label{sec:rq4}

The two remaining comparisons this paper is designed to make --- fine-tuned
GLiNER against the token-classification fleet above, and a zero-shot
frontier LLM against the same fleet, both under identical LLM-gold
evaluation --- require artifacts not yet available: a GLiNER checkpoint
retrained on the current gold (in progress) and a frontier-LLM rerun on
the current gold (not yet performed; the prior Gemini numbers were scored
against dictionary-silver gold and are not reported here to avoid
repeating RQ1's circularity mistake). We report both once available.

\section{Limitations}
\label{sec:limitations}

\paragraph{Training gold is LLM-only, not human-verified.} Human curation
is applied to the 16 validation/test advisories but not to the 66 training
advisories (\S\ref{sec:data-method}), for cost reasons. If Claude Sonnet's
proposals have systematic blind spots of their own --- the same structural
risk that motivated moving off the dictionary annotator in the first
place, just with a different annotator --- training labels would inherit
them even though the evaluation labels, being human-curated, would not.
The results in \S\ref{sec:results} should be read as ``honestly evaluated,''
not as ``free of all annotator bias''; only the measurement side of that
distinction is currently guaranteed.

\paragraph{Blind gold audit not yet complete.} A blind, model-identity-hidden
judgment sheet of 300 predicted spans (\texttt{paper/capstone/gold\_audit.md})
is prepared but not yet human-scored. Once complete, it will report
per-model precision independent of both the dictionary-silver noise that
motivated this pipeline change and any residual bias in the LLM-primary
gold itself, since a human judges each span directly against the source
text rather than against either annotator's labels.

\paragraph{Test set size.} Eight held-out advisories (1{,}995 gold spans in
the test set) is larger than the prior pipeline's test set (1{,}442
dictionary-silver spans) and covers label types dictionary annotation
structurally could not, but the lowest-support labels (\emph{cwe
identifier}, $n=4$; \emph{yara rule name}, $n=9$) still have too few
examples for their per-label numbers to be read as precise.

\paragraph{Windowing artifacts.} All three token classifiers resolve each
word's tag from whichever sliding window contains it most centrally at
inference time, which is a heuristic, not a principled ensembling method;
a small fraction of predictions near window boundaries may be affected.

\paragraph{RQ4 incomplete.} The architecture comparison against GLiNER and
the frontier-LLM comparison both require artifacts not yet produced
(\S\ref{sec:rq4}); until the GLiNER checkpoint is retrained on identical
gold, the smaller circular-to-honest drop it showed in
Table~\ref{tab:circular-vs-honest} cannot be attributed to architecture
rather than to its training data happening to be less circular than the
token classifiers'.

\section{Conclusion}

Evaluating against labels produced by the same annotator being improved
upon is circular, and the gap it hides is not small: a checkpoint that
scored 0.74--0.76 F1 against dictionary-silver gold scores 0.35--0.39
against independently-produced gold (RQ1). Retraining on that same honest
gold recovers most of the apparent quality --- SecureBERT reaches 0.71 F1,
nearly matching its own old circular number --- showing the earlier result
was not fabricated so much as measured against the wrong yardstick (RQ2).
Domain-adaptive pretraining (SecureBERT) recovers more honest-gold quality
than raw parameter count (RoBERTa-large), which notably regresses on the
single highest-volume label rather than improving on it. Structured,
pattern-shaped entity types are essentially solved by all three
architectures; open-vocabulary, context-dependent types --- most visibly
\emph{malware family} and \emph{threat actor group} --- remain the hard
core of this task regardless of training volume (RQ3).

The architecture question this project set out to answer --- whether
GLiNER's span-matching design generalizes better than a token classifier's
softmax head once measured honestly --- is not yet answered. \pending{RQ4
is pending a GLiNER checkpoint retrained on the identical LLM-gold labels
(training in progress) and a frontier-LLM rerun on the same gold; both
will be reported once available, along with the blind gold-audit results
(\S\ref{sec:limitations}).}

\bibliography{capstone}
\bibliographystyle{acl_natbib}

\end{document}
